% !TEX root = ../main.tex
\chapter{Verwandte Arbeiten}
\label{ch:related_work}

Dieses Kapitel ordnet das Projekt \glqq Ascenta\grqq\ in das Feld bestehender technischer Lösungen ein und zeigt die unterschiedlichen konzeptionellen Ansätze auf.

\section{Klassische Mobilitätshilfen}
Der weiße Langstock dient der taktilen Erfassung von Bodenhindernissen und Stufen. \glqq Ascenta\grqq\ ist als komplementäres System konzipiert, das durch den Einsatz von Time-of-Flight-Sensorik zusätzliche Informationen über den Raum oberhalb der Gürtellinie liefert, ohne die Funktion des Langstocks zu ersetzen

\section{Smartphone-basierte Applikationen}
Anwendungen wie \textit{Seeing AI} oder \textit{Be My Eyes} nutzen die im Smartphone integrierte Sensorik zur Umgebungsanalyse. Diese Lösungen erfordern die aktive Handhabung des Mobilgeräts. Technisch basieren diese Systeme häufig auf Cloud-Computing-Modellen. \glqq Ascenta\grqq\ verfolgt hingegen einen Wearable-Ansatz mit lokaler Datenverarbeitung (Edge-AI), wodurch die Interaktion freihändig erfolgt und die Funktionalität unabhängig von der Netzabdeckung gegeben ist.

\section{Kommerzielle Standalone-Systeme}
Systeme wie \textit{OrCam MyEye} oder Lösungen auf Basis von \textit{Google Glass} integrieren Kamera und Recheneinheit in einem dedizierten Gehäuse. Diese Geräte sind als spezialisierte Gesamtsysteme konzipiert.

\section{Positionierung von Ascenta}
Das Konzept von \glqq Ascenta\grqq\ unterscheidet sich durch die konsequente Aufteilung der Komponenten (\textit{Split-Brain}). Anstatt alle Rechenprozesse in der Brille zu bündeln, wird das Smartphone als primäre Logikeinheit genutzt. Durch die Verwendung von Standardkomponenten und die lokale Ausführung der Machine-Learning-Modelle (FOMO und Vosk) wird ein System realisiert, das die Datensouveränität durch den Verzicht auf Cloud-Schnittstellen strukturell im Design verankert.